\chapter{Sneezing}

\section*{Definition}
Forceful expulsion of air through nose/mouth via coordinated reflex to remove irritants.

\section*{What Happens in Your Body}
Common response to allergens viral infections nasal irritation helps clear upper airways.

\section*{Causes}
\begin{itemize}
  \item Allergic rhinitis
  \item Viral infection
  \item Irritants/smoke
  \item Vasomotor rhinitis
  \item Bright light reflex
\end{itemize}

\section*{When to See a Doctor}
See your doctor if:
\begin{itemize}
  \item Seek medical advice if sneezing is persistent, interferes with sleep or daily activities, or occurs with breathing difficulty, facial swelling, high fever, or recurrent nosebleeds.
\end{itemize}

\section*{Related Symptoms}
\begin{itemize}
  \item Runny nose
  \item Congestion
  \item Itchy eyes/nose
  \item Watery eyes
  \item Postnasal drip
\end{itemize}

\section*{Self-Care / Home Management}
\begin{itemize}
  \item Reduce exposure to known allergens and irritants, wash hands, and use saline nasal spray if helpful.
  \item An antihistamine or intranasal corticosteroid may help allergic rhinitis when used according to the label or a clinician’s advice.
\end{itemize}

\section*{How Your Doctor Will Evaluate You}
\begin{itemize}
  \item History and examination assess allergic triggers, infection, medication exposure, nasal obstruction, and associated eye or breathing symptoms.
  \item Allergy testing or an ear, nose, and throat assessment may be considered for persistent or recurrent symptoms.
\end{itemize}

\section*{Treatment Options}
\begin{itemize}
  \item Treatment is directed at the cause, commonly allergic rhinitis, a viral infection, irritant exposure, or less often another nasal disorder; antibiotics are not routine for uncomplicated sneezing.
\end{itemize}

\section*{References}
\begin{thebibliography}{99}
\bibitem{sneezing:ref1} Petit R. ``The sneeze reflex.'' \textit{Otolaryngology--Head and Neck Surgery}, 2013;148(4):535--537.
\bibitem{sneezing:ref2} Bousquet J, Khaltaev N, Cruz AA, et al. ``Allergic Rhinitis and its Impact on Asthma (ARIA) 2008 update.'' \textit{Allergy}, 2008;63(Suppl 86):8--160.
\bibitem{sneezing:ref3} Seidman MD, Gurgel RK, Lin SY, et al. ``Clinical practice guideline: Allergic rhinitis.'' \textit{Otolaryngology--Head and Neck Surgery}, 2015;152(1 Suppl):S1--S43.
\bibitem{sneezing:ref4} Baraniuk JN, Merck SJ. ``Neuroregulation of human nasal mucosa.'' \textit{Annals of the New York Academy of Sciences}, 2009;1170:604--609.
\bibitem{sneezing:ref5} Gwaltney JM Jr. ``Clinical significance and pathogenesis of viral respiratory infections.'' \textit{American Journal of Medicine}, 2002;112(6A):13S--18S.
\bibitem{sneezing:ref6} Eccles R. ``Mechanisms of the common cold: An overview.'' \textit{The Lancet}, 2005;365(9469):50--56.
\end{thebibliography}
